\begin{originalresult}{Theorem 6.1 (Zhang; author source p.36)}
If a two-party gadget $g:X\times Y\to\bits$ satisfies $\AND,\OR\le g$, then for every function $f:\bits^n\to Q$, the function $f\comp g$ has (constant error) randomized communication complexity $\Omega(\operatorname{bs}(f))$.
\end{originalresult}
For completeness, we present the proof that is a reduction from unique set disjointness.
\begin{originalresult}{Definition 6.2}
The disjointness function $\operatorname{Disj}_n:\bits^n\times\bits^n$ is defined as follows:
\[\operatorname{Disj}_n(x,y):=\bigvee_{i=1}^n x_i\wedge y_i.\]
Unique set disjointness, $\operatorname{PromiseDisj}_n$, is a promise version of disjointness with the promise that there is always at most one $i$ such that $x_i\wedge y_i=1$.
\end{originalresult}
\begin{proof}[Proof of Theorem 6.1]
Razborov (1992) showed that $R(\operatorname{PromiseDisj}_m)=\Omega(m)$. We will show a reduction from this problem to computing $f\comp g$. To illustrate the proof, let us start with the case in which $f(0)=0$ and $f(e_i)=1$ for all $i\in[n]$, where $e_i$ is the $i$'th unit vector. Then
\[\operatorname{PromiseDisj}_n(x,y)=f(x_1\wedge y_1,\ldots,x_n\wedge y_n).\]
Since $\AND\le g$, we can solve unique set disjointness using a protocol for $f\comp g$. Therefore, $f\comp g$ has communication complexity $\Omega(n)$. Note that this case only required $\AND\le g$; we need $\OR\le g$ in order to handle inputs changing from 1 to 0.
To handle the general case, let inputs $X^\wedge(0),X^\wedge(1)\in X$ and $Y^\wedge(0),Y^\wedge(1)\in Y$ witness $\AND\le g$, and let $X^\vee(0),X^\vee(1)\in X$ and $Y^\vee(0),Y^\vee(1)\in Y$ witness $\OR\le g$.
Let $z$ be a point witnessing $\operatorname{bs}(f)$, and let $B_1,\ldots,B_m$ be the corresponding blocks, where $m=\operatorname{bs}(f)$. Given inputs $x,y$ to unique set disjointness, we construct the following inputs $X,Y$ to $f\comp g$:
\begin{itemize}
\item If $i$ does not belong to any of the blocks, then we fix $X_i,Y_i$ so that $g(X_i,Y_i)=z_i$.
\item If $i\in B_j$ and $z_i=0$, put $X_i=X^\wedge(x_j),Y_i=Y^\wedge(y_j)$.
\item If $i\in B_j$ and $z_i=1$, put $X_i=X^\vee(\bar x_j),Y_i=Y^\vee(\bar y_j)$.
\end{itemize}
The reader can check that if $x,y$ are disjoint, then applying $g$ to $X,Y$ yields $z$, and if they intersect only at $j$, then it yields $z\oplus B_j$. Hence, $f\comp g$ can be used to compute unique set disjointness of length $m$, and consequently, we obtain a lower bound of $\Omega(m)$.
\end{proof}
